\subsection{R 中的乘法}

有理直线 \(\mathbf{Q}\) 的拓扑不仅与加法群结构相容，而且与 \(\mathbf{Q}\) 的域结构相容。因为函数 \(xy\) 在 \((0, 0) \in \mathbf{Q} \times \mathbf{Q}\) 处连续，这是由于对每个整数 \(n > 0\) ，关系 \(|x| \leqslant 1/n\) 与 \(|y| \leqslant 1/n\) 合起来蕴含 \(|xy| \leqslant 1/n^2 \leqslant 1/n\) ；另一方面，若 \(a\) 是任一非零有理数，则函数 \(ax\) 在 \(x = 0\) 处连续，这是由于对每个整数 \(n > 0\) ，关系 \(|x| \leqslant 1/n |a|\) 蕴含 \(|ax| \leqslant 1/n\) 。这表明 \(xy\) 在 \(\mathbf{Q} \times \mathbf{Q}\) 的每一点连续（第三章 § 6 第 3 目）。

为了证明 \(1 / x\) 在 \(\mathbf{Q}^*\) 上连续，我们更精确地证明：\(1 / x\) 在 0 的任一邻域 \(\mathbf{V}\) 的余集中（关于加法结构）一致连续。即我们有 \(\left|\frac{1}{x} - \frac{1}{y}\right| = \frac{|x - y|}{xy}\) ；存在整数 \(m > 0\) ，使得 \(|x| \geqslant 1 / m\) 对每个 \(x \in \mathbb{C}\mathbf{V}\) 成立；若 \(x\) 与 \(y\) 是 \(\mathbb{C}\mathbf{V}\) 中满足 \(|x - y| \leqslant 1 / m^2 n\) 的任意两点，则将有 \(\left|\frac{1}{x} - \frac{1}{y}\right| \leqslant \frac{1}{n}\) 。

在函数 \(1 / x\) 下，\(\mathbf{Q}^*\) 上（关于加法一致结构）的任一不以 \(0\) 为聚点的 Cauchy 滤子的像是（关于加法一致结构的）一个 Cauchy 滤子。因此（第三章 § 6 第 8 目，命题 7）：

命题 1. 分别定义在 Q × Q 与 Q* 上的函数 xy 与 1/x 可以分别连续延拓到 R × R 与 R* 上，并在 R 上定义一个域结构。赋有这个结构的 R 称为实数域。

第三章 § 6 中建立的拓扑域的所有性质当然都适用；特别地，每个以实数为系数的 n 个实变量的有理函数，在其分母不为零的 \(R^{n}\) 的每一点处连续。

